\section{Remove the wall: nuisance and dynamics interact}
\label{sec:cubes}

Cube worlds reduce the scene to one moving object. Random backgrounds either change independently
each frame or remain fixed for an episode; dynamics are deterministic or include hidden wind.
A fixed background is predictable but irrelevant to position. Per-frame noise is unpredictable and
also irrelevant. These cases distinguish visual complexity from predictive usefulness.

The strongest replicated comparison holds the 512-dimensional grid architecture, dataset size
($500{,}000$ transitions), and training length ($8{,}500$ updates) fixed across four gray-noise
conditions. A supervised nonlinear position read-out is fitted on training encodings and evaluated
on independent evaluation frames. It achieves subpixel error whenever either the background is
persistent or dynamics are deterministic. Combining per-frame noise with wind raises error to
$21.33\pm0.06$ pixels across three model initializations (Table~\ref{tab:cube-main}). This is a
failure to expose position through the tested read-out. Its prediction loss remains high, about
$0.73$ per coordinate, so it is not evidence of a successfully optimized low-loss nuisance solution.

\begin{figure}[t]
\centering
\paperfig{readout_by_seed.pdf}{1.0}
\caption{\textbf{The matched gray-noise cube comparison.} Position read-out from real encodings
for three independently initialized grid JEPAs in each condition. Architecture,
latent dimension, update count, and training dataset size are fixed. A high probe error is a failure
of the tested read-out, not proof of absent information.}
\label{fig:cube-readout}
\end{figure}

\begin{table}[t]
\centering
\caption{\textbf{A matched nuisance--dynamics interaction.} Grid-JEPA position L1 error in pixels,
mean $\pm$ sample standard deviation across three model fits. The same supervised read-out is
applied to real encodings and eight-step predictions; the latter can also reflect read-out
extrapolation. Predicted-window evaluation excludes resets and uses overlapping windows.}
\label{tab:cube-main}
\inputtable{cube_main.tex}
\end{table}

Other fitted models fail at different interfaces. On a black background without wind, the
class-token JEPA permits $0.087$-pixel position read-out from real encodings, but that read-out has
$45.27$-pixel error on eight-step predictions. A flow model also fails on persistent color texture,
with $18.56$-pixel real-latent error. These are single-fit observations with different dataset sizes,
reported separately in Appendix~\ref{app:cubes}; they preclude a uniform ranking of model families.
Together, the cube results motivate an objective-level question rather than a universal empirical
mechanism: what does low prediction loss plus non-collapse actually guarantee?
